\chapter[\emj{🗃️}~Hash Tables (Tablas Hash / Mapas)]{\emj{🗃️}~Hash Tables (Tablas Hash / Mapas)}

\begin{dsaquees}

Un casillero de escuela 🎒. Cada estudiante tiene un número de
casillero asignado por una fórmula mágica (la Hash Function).

Para guardar la mochila de "Carlos", aplicas la fórmula: 
Hash("Carlos") = 42. Vas directo al casillero 42 y la guardas. 
Para buscarla, vuelves a aplicar la fórmula: Hash("Carlos") = 42,
y vas DIRECTO al casillero 42.

¡No tienes que revisar todos los casilleros uno por uno! Es como 
tener un acceso instantáneo a cualquier dato si conoces su "llave".

\end{dsaquees}

\begin{dsaotraforma}

Una Hash Table es una estructura que asocia Llaves (Keys) con Valores
(Values). Usa una función matemática para convertir la Key en un 
índice de un array.

El objetivo principal es lograr operaciones de búsqueda, inserción y 
eliminación en tiempo constante promedio O(1). Es la estructura más 
eficiente para búsquedas rápidas en la industria.

\end{dsaotraforma}

\begin{dsadiagrama}
\begin{verbatim}
Key: "Carlos"  --> [ Hash Function ] --> Index: 2
Key: "Maria"   --> [ Hash Function ] --> Index: 5

Array Interno (Buckets):
[0] : NULL
[1] : NULL
[2] : {Key: "Carlos", Value: "Mochila Azul"}
[3] : NULL
[4] : NULL
[5] : {Key: "Maria",  Value: "Mochila Roja"}

Colisión (Collision):
Si Hash("Juan") también da 2, ocurre una colisión. 
Solución común: Chaining (una Linked List en el índice 2).

[2] : ["Carlos"] -> ["Juan"] -> NULL
\end{verbatim}
\end{dsadiagrama}

\begin{lstlisting}
1. Hash Function: Convierte cualquier Key (string, objeto, etc.) en
   un número entero dentro del rango del array.
2. Insert (set): Calculas el hash, vas al índice y guardas el par {K,V}.
3. Search (get): Calculas el hash, vas al índice y devuelves el valor.
4. Delete: Calculas el hash y eliminas la entrada en ese índice.
5. Rehashing: Si la tabla se llena mucho (Load Factor > 0.75), se crea
   un array más grande y se vuelven a calcular todos los hashes.

📝 PSEUDOCÓDIGO
──────────────────────────────────────────────────────────────

// Insertar en un Hash Map (con Chaining)
función insertar(llave, valor):
    índice = hash(llave) % tamaño_tabla
    lista = buckets[índice]
    SI llave ya existe en lista:
        actualizar valor
    SINO:
        agregar {llave, valor} al final de la lista

💻 CÓDIGO C++
──────────────────────────────────────────────────────────────

#include <iostream>
#include <unordered_map> // Hash Map estándar en C++
#include <unordered_set> // Hash Set (solo llaves únicas)
#include <string>
using namespace std;

int main() {
    // 1. Declaración
    unordered_map<string, int> edades;

    // 2. Insertar -> O(1) promedio
    edades["Carlos"] = 25;
    edades["Maria"] = 30;
    edades.insert({"Juan", 22});

    // 3. Acceder -> O(1) promedio
    cout << "Edad de Maria: " << edades["Maria"] << endl;

    // 4. Verificar existencia
    if (edades.find("Pedro") == edades.end()) {
        cout << "Pedro no está en la tabla." << endl;
    }

    // 5. Eliminar -> O(1)
    edades.erase("Carlos");

    // 6. Recorrer
    for (auto const& [nombre, edad] : edades) {
        cout << nombre << " tiene " << edad << " años." << endl;
    }

    return 0;
}

⏱️ COMPLEJIDAD (Big-O)
──────────────────────────────────────────────────────────────

  Operación   │ Promedio   │ Peor caso
  ────────────┼────────────┼──────────────
  Insert      │ O(1)       │ O(n)
  Search      │ O(1)       │ O(n)
  Delete      │ O(1)       │ O(n)
  Espacio     │ O(n)       │ O(n)

* El peor caso O(n) ocurre si la función hash es muy mala y todos los
  elementos caen en el mismo casillero (colisión masiva).
\end{lstlisting}

\begin{dsaentrevistas}

✅ Regla de Oro: Si necesitas buscar algo en O(1), LA respuesta es
   un Hash Map o Hash Set. 

💡 "Two Sum": El problema más famoso que se resuelve con Hash Map. 
   Guardas el complemento de lo que buscas para encontrarlo en O(1).

⚠️ map vs unordered\_map (C++):
   - `unordered\_map`: Usa Hash Table, O(1) promedio, no tiene orden.
   - `map`: Usa Red-Black Tree, O(log n), mantiene las llaves ordenadas.

🔑 Collision Handling: Debes saber explicar Chaining (listas) y 
   Open Addressing (buscar el siguiente hueco libre).

\end{dsaentrevistas}

\begin{dsaresumen}

• Key -> Hash Function -> Index.
• Acceso, búsqueda e inserción instantáneos O(1).
• La función hash debe ser rápida y distribuir bien los datos.
• Colisiones: Cuando dos llaves quieren el mismo lugar.
• Load Factor: Qué tan llena está la tabla.
• Estructura más usada en la vida real por su velocidad.
════════════════════════════════════════════════════════════════

\end{dsaresumen}


